% Replacement text for the Executive Summary and Abstract.
% The executive summary contains no figures or references.

\section*{Executive Summary}

Sustained lunar surface operations will require engineering environments able to coordinate systems, stakeholders, data sources, and operational assumptions within an evolving mission definition. ECLIPSE was developed as a Digital Engineering Ecosystem for collaborative lunar-base architecting. It connects a SysML-based authoritative source of truth with module selection, site assessment, urban planning, discrete-event simulation (DES), CAD submission, and NVIDIA Omniverse visualization. Users can select systems, place them at a candidate site, generate routes, evaluate operational behavior, associate CAD geometry with the selected systems, and inspect the resulting mission in three dimensions.

The original implementation demonstrated this chain through an in-situ resource utilization (ISRU) scenario. However, several links remained manually coupled or dependent on hardcoded logic. Visualization was prepared around a known scenario rather than generated from a reusable scenario package. CAD submission primarily expected USD-family assets. Most importantly, the DES was implemented as an ISRU proof of concept: although the interface exposed several sliders, resource-flow logic, rover behavior, storage assumptions, thresholds, and power events remained largely embedded in Python. Users could provide an architecture and receive results, but could not fully inspect, modify, or define the operational scenario that produced them.

This work presents an automated, flexible, and coherent extension of ECLIPSE that closes critical gaps between mission definition, simulation, CAD data, and immersive visualization. It transforms previously manual or hardcoded handoffs into explicit, scenario-specific data products, enabling lunar mission scenarios to be configured, simulated, preserved, compared, and visualized through a unified digital engineering workflow. The work addresses three integrated developments: simulation-to-Omniverse integration, multi-format CAD submission with SysML metadata propagation, and a more transparent and extensible DES framework.

First, the simulation-to-Omniverse workflow was reorganized around a scenario-specific visualization manifest. After a run, the pipeline stores the scene USD, waypoint USD, terrain reference, DES results and log, actors, and terrain-projection data with the scenario. The Omniverse extension reads the selected package rather than defining timing, routes, or placement internally. It provides scenario selection, playback controls, route visualization, camera modes, and a telemetry dashboard synchronized with recorded DES time. This is not a direct stream from a concurrently running DES. The Urban Planning map frame is propagated so that terrain, module placement, and rover routes share a metric coordinate system. A terrain sampler grounds module footprints and lifts rover routes to the prepared terrain mesh; local slope can be visualized along the route. Automatic extraction of arbitrary terrain patches remains future work.

Second, CAD submission was extended to STEP/STP, STL, OBJ, and USD-family files. STEP/STP geometry is converted through OpenCascade and pythonOCC, STL and OBJ meshes through trimesh, and USD files through the USD Python API. Each submission produces an Omniverse-compatible USD asset and a browser preview. The workflow preserves source-CAD metadata, such as provenance, surface area, center of mass, and volume when reliable, separately from visualization metadata, such as generated paths, bounding box, units, mesh count, materials, and orientation. SysML remains authoritative for module dimensions: CAD geometry is scaled to the SysML target dimensions, and scale-dependent physical metadata is written into an auto-generated SysML CAD metadata block.

Third, the DES was moved beyond its original role as a fixed ISRU demonstration. Previously embedded assumptions are now collected in a structured configuration exposed through the interface. Users can define module multiplicities, rover counts and payloads, resource routes, transport and processing parameters, storage limits, power generation, battery settings, and module power-demand models. The ISRU engine supports multiple rovers and plants, per-plant input storage, route-derived haul distances, rover return travel, and configurable power generation and consumption. ISRU processing and LOX-storage energy are now passed to the power manager.

An initial generic scenario engine was developed in parallel. It compiles placed modules, SysML interfaces, resource routes, rover assignments, equations, and power configuration into SimPy processes. Modules are classified from their graph, ports, flow directions, and attributes as Sources, Processors, Storage elements, Consumers, Transporters, Power Generators, or Passive elements. Each class has a standard operational behavior and required equations. User equations use a restricted parser and affect the DES only when they map to supported runtime behavior. The framework has been exercised with cargo-relay, ice-processing, and water-transfer examples. It is not yet a universal lunar-operations simulator.

The resulting prototype can save and reload named scenarios, compare selected time histories and final measures of effectiveness, generate scenario-specific Omniverse packages, and visualize each completed scenario independently. Overall, this work moves ECLIPSE from a manually coupled demonstration toward an automated, flexible, and coherent digital engineering workflow for lunar mission design.

\begin{abstract}
ECLIPSE is a Digital Engineering Ecosystem for collaborative lunar-base architecting that connects a SysML-based authoritative source of truth, site selection and urban planning, discrete-event simulation (DES), CAD submission, and NVIDIA Omniverse visualization. This paper presents an automated, flexible, and coherent extension of ECLIPSE that closes critical gaps between mission definition, simulation, CAD data, and immersive visualization. Previously manual or hardcoded handoffs are replaced by explicit, scenario-specific data products that allow lunar mission scenarios to be configured, simulated, preserved, compared, and visualized through a unified workflow.

The contribution combines an automated simulation-to-Omniverse integration based on scenario-specific manifests; a flexible multi-format CAD submission workflow that preserves and propagates engineering metadata to the SysML authoritative source of truth; and a coherent redesign of the discrete-event simulation framework, which makes the original ISRU scenario configurable and establishes a role-based foundation for additional lunar operational scenarios. The extension selects and plays generated packages, provides route and camera tools, and displays telemetry synchronized with recorded simulation time. Terrain sampling and map-frame propagation align prepared terrain meshes, module footprints, and rover routes in a common metric frame. CAD ingestion supports STEP/STP, STL, OBJ, and USD-family files, while source engineering metadata and post-conversion visualization metadata are extracted separately. CAD geometry is scaled to SysML dimensions, and scaled physical metadata is written back into the SysML model. The work demonstrates a more traceable digital thread while identifying remaining limitations in generic behavior modeling, terrain-database automation, and real-time collaborative operation.
\end{abstract}
